% Chapter 6. DEME, the ethics engine. Author: agent D.
% Source conventions for the % src: comments below.
%   lib: = C:\source\erisml-lib at origin/main 2d584e8 (2026-10-03), package erisml-lib 3.1.0.
%          The local checkout of main is 7 commits behind origin/main. The tactical-layer
%          fail_closed option and em_failures record exist only from 0da0b84 (PR #145) on, and
%          the TragicConflictEM rounding from 8af9411 (PR #146) on. Line numbers are origin/main.
%   gtc: = C:\source\gtc-prototype (working tree, 2026-10-03).
\chapter{DEME, the ethics engine}\label{ch:deme}

DEME is the part of the \erisml{} library that judges a set of candidate actions and says which
of them may be taken. It reads a structured description of each candidate, runs a set of ethics
modules over it, removes the candidates that any entitled module forbids, ranks the rest and
writes a hashed record of the whole decision. In the home-care twin of Part~III it is the last
check before the robot's body moves. This chapter describes the engine as the code defines it,
gives the arithmetic of each step, lists every module that exists, and then shows how the twin
wires the engine into its output gate. The chapter ends with a plain account of which parts are
deterministic and which parts rest on a model.
% src: lib:src/erisml/ethics/deme.py:1-3 (name "Democratic Ethics Module Engine")
% src: lib:src/erisml/ethics/layers/pipeline.py:132-322 (DEMEPipeline.decide)
% src: gtc:twin/output_gate.py:1-23

\section{The engine is a three-layer pipeline over structured facts}\label{sec:deme-overview}

The package docstring expands the name as Democratic Ethics Module Engine. The file
\code{deme.py} holds a small class \code{DEME} whose \code{evaluate} method applies only the
deontic maxim gate of Section~\ref{sec:deme-deontic} and otherwise returns \code{ALLOW}. Its own
docstring points to the real engine, the class \code{DEMEPipeline} in
\code{ethics/layers/pipeline.py}. The rest of this chapter is about that class and the parts it
uses.
% src: lib:src/erisml/ethics/deme.py:16-49

The pipeline has three layers. The reflex layer runs cheap veto checks on each option alone. Its
module docstring sets a target of under 100 microseconds and says that a production deployment
would run it in hardware. The tactical layer runs every ethics module on each option that
survived the reflex layer and combines their outputs into one moral vector per option. Its
target is 10 to 100 milliseconds. The strategic layer works over many decisions. In the pipeline
it only stores each finished decision proof in a history list for later analysis.
Figure~\ref{fig:deme-pipeline} shows the six steps of one decision, including the step the twin
adds after the pipeline returns.
% src: lib:src/erisml/ethics/layers/__init__.py:5-15
% src: lib:src/erisml/ethics/layers/reflex.py:5-12
% src: lib:src/erisml/ethics/layers/tactical.py:5-12
% src: lib:src/erisml/ethics/layers/strategic.py:389-396 (record_decision appends the proof)

\begin{figure}[t]
\centering
% Chapter 6. The DEMEPipeline as six numbered panels.
% src: erisml-lib (origin/main 2d584e8) src/erisml/ethics/layers/pipeline.py:132-322
% src: gtc-prototype twin/output_gate.py:179-250 (panel 5)
\begin{tikzpicture}[house,
  pan/.style={minimum width=38mm, minimum height=44mm},
  dtl/.style={detail, text width=34mm, anchor=south}]
% row 1
\node[panel=Purple, pan] (p1) at (0,0) {};
\node[panel=Blue, pan] (p2) at (46mm,0) {};
\node[panel=Green, pan] (p3) at (92mm,0) {};
% row 2 (runs right to left so the flow snakes)
\node[panel=Orange, pan] (p4) at (92mm,-50mm) {};
\node[panel=Red, pan] (p5) at (46mm,-50mm) {};
\node[panel=Teal, pan] (p6) at (0,-50mm) {};
% panel 1
\node[step=Purple] at ([yshift=-5.5mm]p1.north) {1};
\node[stepname=Purple] at ([yshift=-12mm]p1.north) {Facts in};
\node[desc, text width=35mm] at ([yshift=-17.5mm]p1.north) {one EthicalFacts\\per candidate option};
\node[dtl] at ([yshift=2mm]p1.south) {consequences\\rights and duties\\justice, autonomy\\privacy, epistemic\\optional maxim};
% panel 2
\node[step=Blue] at ([yshift=-5.5mm]p2.north) {2};
\node[stepname=Blue] at ([yshift=-12mm]p2.north) {Reflex veto};
\node[desc, text width=35mm] at ([yshift=-17.5mm]p2.north) {cheap checks on\\each option alone};
\node[dtl] at ([yshift=2mm]p2.south) {VetoRule list\\deontic maxim gate\\a rule that raises\\vetoes (fail\_open off)};
% panel 3
\node[step=Green] at ([yshift=-5.5mm]p3.north) {3};
\node[stepname=Green] at ([yshift=-12mm]p3.north) {Tactical vector};
\node[desc, text width=35mm] at ([yshift=-17.5mm]p3.north) {every module judges\\the surviving options};
\node[dtl] at ([yshift=2mm]p3.south) {nine dimensions each\\tier-weighted mean\\veto from tiers 0 to 2\\fail\_closed on error};
% panel 4
\node[step=Orange] at ([yshift=-5.5mm]p4.north) {4};
\node[stepname=Orange] at ([yshift=-12mm]p4.north) {Ranking};
\node[desc, text width=35mm] at ([yshift=-17.5mm]p4.north) {unflagged options by\\scalar collapse};
\node[dtl] at ([yshift=2mm]p4.south) {harm inverted\\weighted mean of nine\\best first\\forbidden list kept};
% panel 5
\node[step=Red] at ([yshift=-5.5mm]p5.north) {5};
\node[stepname=Red] at ([yshift=-12mm]p5.north) {Veto or replace};
\node[desc, text width=35mm] at ([yshift=-17.5mm]p5.north) {the twin's output gate\\acts on the result};
\node[dtl] at ([yshift=2mm]p5.south) {pass if not forbidden\\else best allowed\\else wait and observe\\tragic goes to a human};
% panel 6
\node[step=Teal] at ([yshift=-5.5mm]p6.north) {6};
\node[stepname=Teal] at ([yshift=-12mm]p6.north) {Proof out};
\node[desc, text width=35mm] at ([yshift=-17.5mm]p6.north) {DecisionProof with\\a SHA-256 hash};
\node[dtl] at ([yshift=2mm]p6.south) {facts hash\\layer outputs\\module judgements\\ranked and forbidden};
% flow
\draw[flow] (p1.east) -- (p2.west);
\draw[flow] (p2.east) -- (p3.west);
\draw[flow] (p3.south) -- node[flowlabel, right]{landscape} (p4.north);
\draw[flow] (p4.west) -- (p5.east);
\draw[flow] (p5.west) -- (p6.east);
\end{tikzpicture}

\caption{One DEME decision in six steps. Steps one to four and the proof run inside
\code{DEMEPipeline.decide}, and step five is the twin's output gate acting on the result. Every
step after the first can only remove or reorder options, so no step adds an option the caller
did not offer.}
\label{fig:deme-pipeline}
\end{figure}

The pipeline never sees raw sensor data, text or a model's output. Its only input is a list of
\code{EthicalFacts} records, one per candidate option. The base module protocol states the
contract directly. Ethics modules consume \code{EthicalFacts} and emit judgements, and they
should reason only over those facts, without access to raw domain data, sensors or models. The
facts are therefore the boundary of the engine. Whatever produced them carries the
responsibility for their truth, and Section~\ref{sec:deme-twin} shows how the twin produces
them.
% src: lib:src/erisml/ethics/modules/base.py:1-9, 26-36

\section{EthicalFacts describe one option in fixed blocks}\label{sec:deme-facts}

An \code{EthicalFacts} record names one option by \code{option\_id} and describes it in typed
blocks. Table~\ref{tab:deme-facts} lists the blocks the modules read. Every block except the
option identifier is optional and defaults to \code{None}. Inside each block every field has a
default. Two defaults matter for safety. The field \code{has\_valid\_consent} defaults to true,
and \code{violates\_rights} defaults to false. A producer that omits a field therefore asserts
the permissive value.
% src: lib:src/erisml/ethics/facts.py:36-99 (blocks and defaults), 166-185 (EthicalFacts)

\begin{table}[t]
\centering
\footnotesize
\begin{tabularx}{\linewidth}{l X}
\toprule
Block & Fields read by the modules in this chapter \\
\midrule
\code{consequences} & \code{expected\_benefit}, \code{expected\_harm}, \code{urgency}, \code{affected\_count} \\
\code{rights\_and\_duties} & \code{violates\_rights}, \code{has\_valid\_consent}, \code{violates\_explicit\_rule}, \code{role\_duty\_conflict} \\
\code{justice\_and\_fairness} & \code{discriminates\_on\_protected\_attr}, \code{exploits\_vulnerable\_population}, \code{exacerbates\_power\_imbalance}, \code{prioritizes\_most\_disadvantaged} \\
\code{autonomy\_and\_agency} & \code{has\_meaningful\_choice}, \code{coercion\_or\_undue\_influence}, \code{manipulative\_design\_present}, \code{can\_withdraw\_without\_penalty} \\
\code{privacy\_and\_data} & \code{privacy\_invasion\_level}, \code{secondary\_use\_without\_consent}, \code{data\_minimization\_respected} \\
\code{societal\_and\_environmental} & \code{environmental\_harm}, \code{long\_term\_societal\_risk}, \code{burden\_on\_vulnerable\_groups} \\
\code{virtue\_and\_care} & \code{expresses\_compassion}, \code{respects\_person\_as\_end}, \code{betrays\_trust} \\
\code{procedural\_and\_legitimacy} & \code{followed\_approved\_procedure}, \code{stakeholders\_consulted}, \code{decision\_explainable\_to\_public} \\
\code{epistemic\_status} & \code{uncertainty\_level}, \code{evidence\_quality}, \code{novel\_situation\_flag} \\
\code{maxim} & \code{action\_kind}, \code{polarity} \\
\bottomrule
\end{tabularx}
\caption{The blocks of an \code{EthicalFacts} record and the fields the modules and the vector
mapping read. The record also carries free \code{metadata}, \code{tags} and \code{extra}
dictionaries that no module in this chapter reads.}
\label{tab:deme-facts}
\end{table}
% src: lib:src/erisml/ethics/facts.py:9-185
% src: lib:src/erisml/ethics/moral_vector.py:399-532 (fields read by from_ethical_facts)

The last block, \code{Maxim}, mirrors a structure that the \erisml{} compiler extracts from text.
It holds an action kind and a polarity, affirmed or negated. A negated maxim such as ``did not
promise'' is the maxim of not performing the action.
% src: lib:src/erisml/ethics/facts.py:151-163

The file \code{facts\_v3.py} extends the record to many parties. \code{EthicalFactsV3} keeps one
entry per party in each block (\code{PartyConsequences}, \code{PartyRights} and six more) and
computes distribution statistics over them. One of these is the Gini coefficient. For party
values sorted so that $x_{(1)} \le \dots \le x_{(n)}$ the code computes
\begin{equation}
G = \frac{2\sum_{i=1}^{n} i\, x_{(i)} - (n+1)\sum_{i=1}^{n} x_{(i)}}{n \sum_{i=1}^{n} x_{(i)}},
\label{eq:deme-gini}
\end{equation}
and returns $0$ for fewer than two values or an all-zero list. A V3 record converts to a V2
record and to a rank-2 moral tensor. The collapse function \code{collapse\_facts\_v3\_to\_v2}
accepts a \code{worst\_case} strategy, but the code returns the same aggregate as the default
strategy, and a comment calls this a simplified implementation. A reader should not assume a
pessimistic collapse exists at the facts level.
% src: lib:src/erisml/ethics/facts_v3.py:5-18, 49-70 (Gini), 79-123 (party blocks)
% src: lib:src/erisml/ethics/facts_v3.py:971-1010, 1053, 1155 (EthicalFactsV3, to_v2, to_moral_tensor)
% src: lib:src/erisml/ethics/facts_v3.py:1319-1348 (worst_case returns facts.to_v2())

\section{The moral vector has nine dimensions in a three by three grid}\label{sec:deme-vector}

Every module in DEME version 2 returns a \code{MoralVector}. It has nine named dimensions, each a
number in $[0,1]$. The constructor rejects any value outside that interval. For eight of the
dimensions one means good. For \code{physical\_harm} one means severe harm. The vector also
carries a list of veto flags, a list of reason codes and an optional dictionary of extension
dimensions.
% src: lib:src/erisml/ethics/moral_vector.py:46-136

The nine dimensions come from a three by three matrix. One axis is scope (individual, relational,
collective). The other axis is mode (what matters, who decides, what we know).
Figure~\ref{fig:deme-vector} places each dimension in its cell with its index $k$ in the tensor
order and its default weight in the scalar collapse below. The canonical names are defined once
in the compiler package \code{erisml\_compiler.ir.v3.dimensions}. \code{moral\_tensor.py}
imports them and keeps a literal copy only as a fallback when the compiler is not installed. A
consistency test in the compiler guards the copy.
% src: lib:docs/moralvector_v2_architecture.md:27-44 (the 3x3 matrix)
% src: lib:docs/moralvector_reference.md:8-29
% src: lib:src/erisml/ethics/moral_tensor.py:16-25, 45-77 (single source of truth and fallback)

\begin{figure}[t]
\centering
% Chapter 6. The nine MoralVector dimensions as the 3 by 3 scope-by-mode matrix, with index k and
% the default weight of the scalar collapse.
% src: erisml-lib docs/moralvector_v2_architecture.md:33-37 (matrix)
% src: erisml-lib src/erisml/ethics/moral_tensor.py:66-76 (k order)
% src: erisml-lib src/erisml/ethics/moral_vector.py:31-43 (DEFAULT_DIMENSION_WEIGHTS)
\begin{tikzpicture}[house,
  cell/.style={minimum width=44mm, minimum height=15mm},
  hd/.style={font=\sffamily\bfseries\footnotesize, text=hsTextMid},
  kk/.style={step=#1, minimum size=5.5mm, font=\sffamily\bfseries\scriptsize}]
\node[hd] at (0,12mm) {What matters};
\node[hd] at (46mm,12mm) {Who decides};
\node[hd] at (92mm,12mm) {What we know};
\node[hd, anchor=east] at (-24mm,0) {Individual};
\node[hd, anchor=east] at (-24mm,-18mm) {Relational};
\node[hd, anchor=east] at (-24mm,-36mm) {Collective};
\foreach \x/\y/\c/\k/\n/\w in {
  0/0/Purple/3/autonomy\_respect/1.0,
  46/0/Blue/1/rights\_respect/1.0,
  92/0/Teal/4/privacy\_protection/1.0,
  0/-18/Green/6/virtue\_care/0.7,
  46/-18/Red/0/physical\_harm/{1.0, inverted},
  92/-18/Amber/8/epistemic\_quality/0.5,
  0/-36/Orange/2/fairness\_equity/1.0,
  46/-36/Grey/7/legitimacy\_trust/1.0,
  92/-36/Teal/5/societal\_environmental/0.8}{
  \node[panel=\c, cell] (c\k) at (\x mm,\y mm) {};
  \node[kk=\c] at ([xshift=4.5mm, yshift=-4.5mm]c\k.north west) {\k};
  \node[font=\ttfamily\scriptsize, text=hs\c] at ([xshift=3mm, yshift=1.5mm]c\k.center) {\n};
  \node[desc] at ([xshift=3mm, yshift=-4mm]c\k.center) {weight \w};
}
\end{tikzpicture}

\caption{The nine moral vector dimensions placed by scope and mode, each with its tensor index
and its default weight in the scalar collapse. Only \code{physical\_harm} counts upward as worse,
so the collapse inverts it before weighting.}
\label{fig:deme-vector}
\end{figure}

The reference document for the vector describes each tensor cell as a signed valence in
$[-1,+1]$ with confidence and source metadata. That is the compiler's \code{DimensionScore}. The
library's \code{MoralVector} uses $[0,1]$ and validates it. The two scales differ, and a value
passed from one to the other needs a mapping.
% src: lib:docs/moralvector_reference.md:31-33
% src: lib:src/erisml/ethics/moral_vector.py:122-136 (bounds check [0,1])

\paragraph{Scalar collapse.} Ranking needs one number per option. For a vector $v$ with
components $v_d$ and weights $\alpha_d$, \code{to\_scalar} computes
\begin{equation}
s(v) = \frac{\sum_{d} \alpha_d\, \tilde v_d}{\sum_{d} \alpha_d},
\qquad
\tilde v_d =
\begin{cases}
1 - v_d & d = \code{physical\_harm},\\
v_d & \text{otherwise.}
\end{cases}
\label{eq:deme-scalar}
\end{equation}
The default weights are $1.0$ for harm, rights, fairness, autonomy, privacy and legitimacy,
$0.8$ for societal and environmental effects, $0.7$ for virtue and care, and $0.5$ for epistemic
quality. Extension dimensions enter only when the weight dictionary names them. If all weights
are zero the function returns $0.5$.
% src: lib:src/erisml/ethics/moral_vector.py:29-43 (DEFAULT_DIMENSION_WEIGHTS), 161-203 (to_scalar)

\paragraph{Dominance.} A vector $u$ Pareto-dominates $v$ when it is no worse on every dimension
and strictly better on at least one, with lower harm counted as better. The method
\code{dominates} implements this test, and \code{MoralLandscape.pareto\_frontier} uses it. The
pipeline does not use the frontier. It ranks by the scalar of Equation~\eqref{eq:deme-scalar}.
% src: lib:src/erisml/ethics/moral_vector.py:209-240
% src: lib:src/erisml/ethics/moral_landscape.py:136 (pareto_frontier), 268-287 (rank_by_scalar)

\paragraph{The standard mapping from facts.} \code{MoralVector.from\_ethical\_facts} gives a
default reading of a facts record. A rights violation sets rights respect to $0$ and adds the
flag \code{RIGHTS\_VIOLATION}. Discrimination on a protected attribute sets fairness to $0$ and
adds \code{DISCRIMINATION}. Missing consent without a violation sets rights respect to $0.5$.
Coercion sets autonomy to $0.2$. Privacy is $1 - \code{privacy\_invasion\_level}$, halved for
secondary use without consent. Epistemic quality is $1 - \code{uncertainty\_level}$, capped at
$0.3$ for low evidence and $0.6$ for medium evidence. The tactical layer uses this mapping only
when it is disabled.
% src: lib:src/erisml/ethics/moral_vector.py:398-532
% src: lib:src/erisml/ethics/layers/tactical.py:142-160 (disabled layer returns from_ethical_facts)

\paragraph{Measured dimensions.} The reference document reports learned encoders, called xBSE
feeders, that score single dimensions from text. It reports cross-dataset AUROC between $0.622$
and $0.853$ for seven of the nine dimensions and none yet for \code{rights\_respect} and
\code{societal\_environmental}. The decision proof has a record type that binds such a feeder's
checkpoint hash and pre-registered bar into the proof. No feeder runs in the twin. Every vector
in the twin is computed by the rules below.
% src: lib:docs/moralvector_reference.md:48-62 (feeder table)
% src: lib:src/erisml/ethics/decision_proof.py:80-122 (FeederValidationRecord)

\section{Ethics modules are sorted into five tiers}\label{sec:deme-modules}

An ethics module (EM) is any object with a name, a stakeholder and a \code{judge} method. The
version 1 protocol returns an \code{EthicalJudgement} with a verdict and one normative score in
$[0,1]$. The version 2 protocol, \code{EthicsModuleV2}, adds an integer tier and returns an
\code{EthicalJudgementV2} whose score is a \code{MoralVector}. A verdict is one of
\code{strongly\_prefer}, \code{prefer}, \code{neutral}, \code{avoid} and \code{forbid}. A
version 2 judgement counts as a veto when its verdict is \code{forbid} or its vector carries a
veto flag.
% src: lib:src/erisml/ethics/modules/base.py:26-55 (V1), 154-198 (V2), 277-309 (veto_triggered)
% src: lib:src/erisml/ethics/judgement.py:19-20 (Verdict), 23-89 (V1), 118-182 (V2)

The tier sets both the module's weight in aggregation and whether its veto counts.
Figure~\ref{fig:deme-tiers} gives the defaults of the tactical layer and the modules that exist
at each tier. The tier packages state their intent. Tier 0 modules encode constraints that
governance profiles cannot disable. Tier 1 is for physical safety, tier 2 for rights and
fairness, tier 3 for soft values and tier 4 for checks on the governance system itself. The
packages for tiers 1, 3 and 4 contain no modules. Each holds a comment that names planned modules
(\code{CorePhysicalSafetyEM}, \code{ProxemicsSafetyEM}, \code{BeneficenceEM},
\code{EnvironmentalImpactEM}, \code{PatternGuardEM}, \code{ProfileIntegrityEM}) and an empty
export list.
% src: lib:src/erisml/ethics/layers/tactical.py:46-66 (tier_weights, tier_veto_enabled)
% src: lib:src/erisml/ethics/modules/tier0/__init__.py:5-17
% src: lib:src/erisml/ethics/modules/tier1/__init__.py:5-24, tier3/__init__.py:5-24, tier4/__init__.py:5-24

\begin{figure}[t]
\centering
% Chapter 6. The five module tiers, their default weights and veto rights, and what each holds.
% src: erisml-lib src/erisml/ethics/layers/tactical.py:46-63 (origin/main) weights and veto flags
% src: erisml-lib src/erisml/ethics/modules/tier0..tier4/__init__.py (contents)
% src: gtc-prototype twin/output_gate.py:58-62 (the twin's three modules)
\begin{tikzpicture}[house,
  tier/.style={minimum width=140mm, minimum height=11mm},
  nm/.style={font=\sffamily\bfseries\footnotesize, anchor=west},
  wt/.style={detail, text width=22mm},
  mods/.style={desc, anchor=west, align=left, text width=58mm}]
\foreach \c/\y/\n/\name/\w/\v/\m in {
  Red/0/0/Constitutional/{weight 10.0}/{can veto}/{GenevaEMV2 (twin)\\GenevaEMV3},
  Orange/-13/1/Core safety/{weight 5.0}/{can veto}/{none yet (placeholder)},
  Purple/-26/2/Rights and fairness/{weight 2.0}/{can veto}/{AutonomyConsentEMV2 (twin)\\TriageEMV3, RightsFirstEMV3},
  Blue/-39/3/Soft values/{weight 1.0}/{advisory}/{none in the package\\TragicConflictEM placed here by the twin},
  Grey/-52/4/Meta-governance/{weight 0.5}/{advisory}/{none yet (placeholder)}}{
  \node[panel=\c, tier] (t\n) at (0,\y mm) {};
  \node[step=\c] at ([xshift=6mm]t\n.west) {\n};
  \node[nm, text=hs\c] at ([xshift=11mm]t\n.west) {\name};
  \node[wt] at ([xshift=64mm]t\n.west) {\w\\\v};
  \node[mods] at ([xshift=80mm]t\n.west) {\m};
}
\draw[flow] ([xshift=-4mm]t4.south west) -- node[flowlabel, left, align=right]{more\\weight} ([xshift=-4mm]t0.north west);
\end{tikzpicture}

\caption{The five module tiers with their default aggregation weights, whether their vetoes
count, and the modules that exist at each tier. Tiers 1, 3 and 4 hold no modules of their own,
and the twin places the advisory tragic-conflict module at tier 3.}
\label{fig:deme-tiers}
\end{figure}

\subsection{GenevaEMV2 holds the hard constraints at tier 0}

\code{GenevaEMV2} is registered at tier 0 with weight $10.0$ and veto capability. It starts from
a vector of ones with harm equal to the facts' expected harm. Three facts are hard vetoes. A
rights violation sets $\code{rights\_respect}=0$ and the flag \code{RIGHTS\_VIOLATION}.
Discrimination on a protected attribute sets $\code{fairness\_equity}=0$ and the flag
\code{DISCRIMINATION}. An explicit rule violation sets $\code{legitimacy\_trust}=0$ and the flag
\code{RULE\_VIOLATION}. Four other facts are penalties without a veto. Writing $[\,\cdot\,]$ for
an indicator that is $1$ when its condition holds, the module computes
\begin{align}
\code{autonomy\_respect} &= 0.5^{[\neg\,\text{consent}]},\\
\code{legitimacy\_trust} &= [\neg\,\text{rule violation}]\cdot 0.7^{[\text{role duty conflict}]},\\
\code{fairness\_equity} &= [\neg\,\text{discrimination}]\cdot 0.5^{[\text{exploits vulnerable}]}\cdot 0.8^{[\text{power imbalance}]},\\
\code{epistemic\_quality} &= (1-u)\cdot q,\qquad q \in \{0.5, 0.8, 1\}\ \text{for low, medium, other evidence},
\end{align}
where $u$ is the uncertainty level. Without an epistemic block the epistemic quality is $1$. The
consent penalty applies only when the field
\code{strict\_consent} is true, which is its default. The verdict is \code{forbid} when any flag
is set. Otherwise it is \code{strongly\_prefer} or \code{prefer} when rights, fairness and
legitimacy all exceed $0.7$ (the first when harm is below $0.2$), \code{neutral} when they all
exceed $0.5$, and \code{avoid} otherwise.
% src: lib:src/erisml/ethics/modules/tier0/geneva_em.py:27-32 (registration), 45-56
% src: lib:src/erisml/ethics/modules/tier0/geneva_em.py:73-186 (evaluate_vector)

The module also defines a \code{reflex\_check} that returns true on the same three hard
conditions. No layer of the pipeline calls it. The pipeline's reflex layer runs its own rule
list, described in Section~\ref{sec:deme-pipeline}.
% src: lib:src/erisml/ethics/modules/tier0/geneva_em.py:58-71
% src: lib:src/erisml/ethics/layers/ (no call site of reflex_check; checked with git grep on origin/main)

Matthias argued that an ethical governor built on legal codes cannot express moral
dissent~\cite{matthias2011governor}. \code{GenevaEMV2} is a constraint set of that kind. Its name
refers to the Geneva conventions, and its content is the three hard flags and four penalties
above.

\subsection{AutonomyConsentEMV2 multiplies autonomy penalties at tier 2}

\code{AutonomyConsentEMV2} is registered at tier 2 with weight $2.0$ and veto capability. It
computes one autonomy score as a product of penalties,
\begin{equation}
a = 0.2^{[\text{coercion}]}\cdot 0.4^{[\text{manipulative design}]}\cdot 0.6^{[\neg\,\text{choice}]}\cdot 0.8^{[\neg\,\text{withdrawal}]}\cdot 0.5^{[\neg\,\text{consent}]},
\label{eq:deme-autonomy}
\end{equation}
where the first four factors apply only when an \code{autonomy\_and\_agency} block is present. If
$a < 0.3$ it adds the flag \code{AUTONOMY\_VIOLATION} and the verdict is \code{forbid}. Coercion
with missing consent gives $a = 0.1$ and so a veto. Coercion alone gives $0.2$, also a veto. The
module sets six dimensions to a fixed $0.8$ (fairness, privacy, societal and environmental
effects, virtue and care, legitimacy, and epistemic quality), sets rights respect to $1$ or $0.5$
by consent, and copies harm from the facts.
% src: lib:src/erisml/ethics/modules/tier2/autonomy_consent_em.py:26-31, 34-53, 56-133

\subsection{TragicConflictEM measures how much every option costs}

\code{TragicConflictEM} is a version 1 module. Its docstring says it measures how tragic a choice
is and never vetoes. It sums fixed weights over triggers read from the facts,
\begin{equation}
\begin{aligned}
\tau = \operatorname{round}_6\Bigl(\min\bigl(1,\ &0.20[u_r \ge 0.75] + h_w + 0.15[b \ge 0.6 \wedge h \ge 0.6]\\
&+ 0.25[R] + 0.15[L] + 0.10[\neg C] + 0.15[D]\bigr)\Bigr),
\end{aligned}
\label{eq:deme-tragic}
\end{equation}
where $u_r$ is urgency, $b$ expected benefit, $h$ expected harm, $h_w$ is $0.35$ when
$h \ge 0.8$, $0.25$ when $0.6 \le h < 0.8$ and $0$ otherwise, and $R$, $L$, $C$, $D$ stand for a
rights violation, a rule violation, valid consent and discrimination. The verdict is
\code{neutral} when $\tau \ge 0.55$ and \code{prefer} otherwise, the normative score is
$\max(0,\, 0.85 - 0.6\,\tau)$, and the metadata carries $\tau$, the trigger names and the flag
\code{tragic\_conflict\_high} $= [\tau \ge 0.55]$.
% src: lib:src/erisml/ethics/modules/greek_tragedy_tragic_conflict_em.py:54-142

The rounding to six decimals was added on 2026-10-02. The weights are decimal hundredths, and
summed in binary floating point they drift. The commit message gives $0.20+0.15+0.10$, which
evaluates to $0.44999999999999996$. Near the threshold of $0.55$ such drift would decide the flag
by rounding error.
% src: lib commit 8af9411 "TragicConflictEM: the index is the decimal sum of its weights"
% src: lib:src/erisml/ethics/modules/greek_tragedy_tragic_conflict_em.py:111-115 (comment and round)

\subsection{Adapters let version 1 modules join a version 2 pipeline}

\code{V1ToV2Adapter} wraps a version 1 module and assigns it a tier, $2$ by default. It converts
each judgement with \code{judgement\_v1\_to\_v2}. A score $\sigma$ becomes a vector with every
value dimension equal to $\sigma$, harm equal to $1-\sigma$ and epistemic quality $0.8$. A
\code{forbid} verdict becomes harm $1$, rights respect $0$, the flag \code{V1\_FORBID} and a veto.
The converted judgement carries confidence $0.8$ and a \code{\_migrated\_from\_v1} mark.
\code{V2ToV1Adapter} goes the other way and collapses the vector with the weights of
\code{DEFAULT\_V2\_WEIGHTS}. A version 1 module added to a version 2 pipeline without the adapter
fails, because it has no tier. Before 2026-10-02 the tactical layer dropped that failure
silently.
% src: lib:src/erisml/ethics/modules/base.py:322-386 (adapters)
% src: lib:src/erisml/ethics/judgement.py:107-115, 200-279
% src: lib commit 0da0b84 message; lib:tests/test_tactical_em_failures.py:57-63

One consequence follows for the twin. Wrapped this way, \code{TragicConflictEM} reports
$\sigma = 0.85 - 0.6\,\tau$, so even with $\tau = 0$ its converted vector has harm
$0.15$. This vector enters the weighted mean of Equation~\eqref{eq:deme-aggregate} at the tier 3
weight.

\subsection{The other modules}

Table~\ref{tab:deme-modules} lists every ethics module in \code{ethics/modules} and what each
does. There are two registries. The dictionary \code{EM\_REGISTRY} in the package
\code{\_\_init\_\_} maps four identifiers to version 1 classes. The class \code{EMRegistry}
collects version 2 and version 3 classes through a decorator that records tier, default weight
and veto capability.
% src: lib:src/erisml/ethics/modules/__init__.py:5-31
% src: lib:src/erisml/ethics/modules/registry.py:1-60

\begin{table}[t]
\centering
\footnotesize
\begin{tabularx}{\linewidth}{l c X}
\toprule
Module & Tier & What it does \\
\midrule
\code{GenevaEMV2} & 0 & three hard vetoes, four penalties (above) \\
\code{GenevaEMV3} & 0 & the same constraints per party over a moral tensor, with per-party veto locations \\
\code{AutonomyConsentEMV2} & 2 & product of autonomy penalties, veto below $0.3$ (above) \\
\code{TriageEMV3} & 2 & per-party triage score with vulnerability weights, Gini and maximin metrics \\
\code{RightsFirstEMV3} & 2 & per-party version of \code{RightsFirstEM} \\
\code{CaseStudy1TriageEM} & V1 & forbids rights or rule violations, else a weighted score of benefit $0.35$, harm $0.25$, urgency $0.20$, priority for the disadvantaged $0.15$ and procedure $0.05$, with an epistemic penalty \\
\code{RightsFirstEM} & V1 & forbids rights or rule violations, else \code{prefer} at $0.8$ \\
\code{GenevaBaseEM}, \code{GenevaBaselineEM} & V1 & Geneva scoring with configurable penalty weights (\code{GenevaWeights}) \\
\code{TragicConflictEM} & V1 & tragic-conflict index, never vetoes (above) \\
\bottomrule
\end{tabularx}
\caption{Every ethics module in the library's \code{modules} package. V1 marks a version 1
module, which has no tier until an adapter gives it one.}
\label{tab:deme-modules}
\end{table}
% src: lib:src/erisml/ethics/modules/tier0/geneva_em_v3.py:1-16, 39-69
% src: lib:src/erisml/ethics/modules/triage_em_v3.py:1-16, 40-72, 506-526
% src: lib:src/erisml/ethics/modules/triage_em.py:11-31 (weights), 164-210
% src: lib:src/erisml/ethics/modules/geneva_base_em.py:1-11, 27-45, 72-79, 149-157
% src: lib:src/erisml/ethics/modules/base_v3.py:1-12, 303-325, 372-396 (V3 interfaces and adapters)

\section{The pipeline vetoes first and ranks what remains}\label{sec:deme-pipeline}

\subsection{The reflex layer}

The reflex layer holds a list of \code{VetoRule} objects. Each rule is a name, a category and a
predicate on the facts. The layer checks the rules in order and returns at the first that fires.
If a rule raises an exception the option is vetoed, unless the configuration sets
\code{fail\_open}, which defaults to false. A failure of the whole layer is handled the same way.
% src: lib:src/erisml/ethics/layers/reflex.py:26-104 (categories, VetoRule, config), 150-207 (check)

The layer has three default rules, for rights violation, discrimination and explicit rule
violation. They are installed only when the layer is built with no configuration.
\code{PipelineConfig} always passes a fresh \code{ReflexLayerConfig}, whose rule list is empty.
A \code{DEMEPipeline} built with the default configuration therefore runs no reflex rules. We
checked this by constructing both objects, and the pipeline's reflex layer held zero rules while
a bare \code{ReflexLayer()} held three. In such a pipeline the three hard vetoes still take
effect, but in the tactical layer, through \code{GenevaEMV2}.
% src: lib:src/erisml/ethics/layers/reflex.py:114-148 (defaults only when config is None)
% src: lib:src/erisml/ethics/layers/pipeline.py:44-50 (PipelineConfig default_factory)
% src: checked 2026-10-03 with python, len(DEMEPipeline().reflex.config.veto_rules) == 0

In the same phase the pipeline applies the deontic maxim gate to each option. An option whose
maxim fails the gate joins the reflex-vetoed set, and the reason is recorded.
% src: lib:src/erisml/ethics/layers/pipeline.py:155-170

\subsection{The deontic maxim gate}\label{sec:deme-deontic}

The gate in \code{deontic\_gate.py} applies a universalizability test in the manner of
Kant~\cite{kant1998groundwork} to an action kind and its polarity. It mirrors a knowledge base in
the compiler and copies it so that the library does not depend on the compiler. The knowledge
base has the three sets of Table~\ref{tab:deme-deontic}. The gate vetoes an affirmed prohibition
(a contradiction in conception) and a negated imperfect duty (a contradiction in will). It passes
every other known case. An absent maxim or an unknown action kind never vetoes.
% src: lib:src/erisml/ethics/deontic_gate.py:1-136

\begin{table}[t]
\centering
\footnotesize
\begin{tabularx}{\linewidth}{l X l l}
\toprule
Set & Action kinds & Affirmed & Negated \\
\midrule
Prohibitions & \code{deceive}, \code{break\_commitment}, \code{cheat}, \code{impose\_externality}, \code{coerce}, \code{coerce\_or\_be\_coerced}, \code{inflict\_harm}, \code{use\_as\_means}, \code{refuse} & veto & pass \\
Imperfect duties & \code{help}, \code{protect} & pass & veto \\
Permissible & \code{make\_or\_keep\_commitment}, \code{protect}, \code{help}, \code{disclose}, \code{act\_under\_norm}, \code{act\_under\_authority} & pass & pass \\
\bottomrule
\end{tabularx}
\caption{The deontic gate's knowledge base and its result for each polarity. The kinds
\code{help} and \code{protect} appear in two sets, and the imperfect-duty rule decides their
negated case.}
\label{tab:deme-deontic}
\end{table}
% src: lib:src/erisml/ethics/deontic_gate.py:29-55, 101-136

The kind \code{refuse} is a prohibition, so an affirmed maxim of refusing is vetoed. This is the
robot's maxim, not the person's. The twin's facts carry no maxim, so in the twin the gate never
fires.
% src: lib:src/erisml/ethics/deontic_gate.py:40
% src: gtc:twin/output_gate.py:135-168 (EthicalFacts built with no maxim)

\subsection{The tactical layer}

Every option that survives the reflex layer goes to every module. Let $J(o)$ be the judgements
returned for option $o$, $t(j)$ the tier of the module behind judgement $j$, $w_t$ the tier
weight and $v_{j,d}$ dimension $d$ of its vector. The tactical layer combines them by a
normalized weighted mean,
\begin{equation}
\bar v_d(o) = \frac{\sum_{j \in J(o)} w_{t(j)}\, v_{j,d}}{\sum_{j \in J(o)} w_{t(j)}},
\qquad
(w_0,\dots,w_4) = (10,\ 5,\ 2,\ 1,\ 0.5),
\label{eq:deme-aggregate}
\end{equation}
with weight $1.0$ for any tier outside the table. The veto flags and reason codes of all
judgements are joined without duplicates. With no judgements the result is the default vector.
% src: lib:src/erisml/ethics/layers/tactical.py:46-56 (weights), 233-283 (_aggregate_vectors)

An option is tactically vetoed when some module's judgement triggers a veto and that module's
tier is allowed to veto. Tiers 0 to 2 may, tiers 3 and 4 may not. A tier missing from the table
may veto.
% src: lib:src/erisml/ethics/layers/tactical.py:57-66, 188-195

A module that raises an exception has not judged the option. Since commit 0da0b84 the layer
records each failure as a string with the module name and the exception, logs it, and continues
with the other modules. With \code{fail\_closed} set, the failure also vetoes the option, as an
objection would. The default is false, which keeps the earlier behaviour of judging the option
by the modules that answered. The configuration's docstring says that safety-critical
deployments should set it to true. The pipeline reports the failures per option in
\code{DecisionResult.em\_failures} and in the tactical layer's entry in the proof.
% src: lib:src/erisml/ethics/layers/tactical.py:68-76 (fail_closed), 100-101, 160-208
% src: lib:src/erisml/ethics/layers/pipeline.py:94-95, 191, 207-208, 239, 322
% src: lib:tests/test_tactical_em_failures.py:43-91

\subsection{Ranking and selection}

The forbidden set is the union of the reflex-vetoed and tactically vetoed options. Ranking runs
over a separate set, the options in the moral landscape whose aggregated vector carries no veto
flag. The pipeline sorts these by the scalar $s(\bar v(o))$ of
Equation~\eqref{eq:deme-scalar}, best first. It selects the top option, except that a supplied
baseline option that is not forbidden is selected instead when
\begin{equation}
s(\bar v(\text{baseline})) \ \ge\ 0.95\ s(\bar v(\text{top})).
\label{eq:deme-baseline}
\end{equation}
% src: lib:src/erisml/ethics/layers/pipeline.py:245-264
% src: lib:src/erisml/ethics/moral_landscape.py:180-199 (filter_vetoed uses has_veto), 268-287

The two sets need not agree. A veto from a failing module under \code{fail\_closed} adds no veto
flag to the vector, and the top-ranked option is chosen without a check against the forbidden
set. So the pipeline's \code{selected\_option\_id} can name an option that is also in
\code{forbidden\_options}. A caller should filter the ranking against the forbidden list, as the
twin does, rather than trust the selected option alone.
% src: lib:src/erisml/ethics/layers/pipeline.py:245-256 (selected_id = ranked[0][0], no forbidden check)
% src: gtc:twin/output_gate.py:240-247 (the twin filters ranked against forbidden and allowed)

\section{Each decision leaves a hashed proof}\label{sec:deme-proof}

When \code{generate\_proofs} is set, which is the default, the pipeline builds a
\code{DecisionProof} and finalizes it. Figure~\ref{fig:deme-proof} groups its fields. The input
group holds a hash of the facts, the profile name and hash, the module catalogue version and the
names of the active modules. The facts hash is computed in two steps. Each option's facts are
serialized as sorted-key JSON and hashed with SHA-256, and the list of those hashes is hashed
again. The intermediate group holds the reflex and tactical layer outputs (timings, vetoed
options, deontic reasons, module failures) and one record per module judgement with a SHA-256
hash of its vector. The output group holds the candidates, the selected option, the ranking, the
forbidden options, the rationale and the aggregated vector of each option.
% src: lib:src/erisml/ethics/layers/pipeline.py:284-310
% src: lib:src/erisml/ethics/decision_proof.py:31-78, 125-201, 337-383

\begin{figure}[t]
\centering
% Chapter 6. The DecisionProof: field groups, the canonical hash, and the chain link.
% src: erisml-lib src/erisml/ethics/decision_proof.py:125-271 (fields, compute_hash, canonical dict)
% src: erisml-lib src/erisml/ethics/decision_proof.py:386-453 (DecisionProofChain)
\begin{tikzpicture}[house,
  grp/.style={minimum width=36mm, minimum height=34mm},
  dtl/.style={detail, text width=31mm, anchor=north}]
\node[panel=Purple, grp] (in) at (0,0) {};
\node[panel=Green, grp] (mid) at (40mm,0) {};
\node[panel=Orange, grp] (out) at (80mm,0) {};
\node[stepname=Purple] at ([yshift=-4mm]in.north) {Input state};
\node[dtl] at ([yshift=-8mm]in.north) {input\_facts\_hash\\profile name and hash\\EM catalogue version\\active module names};
\node[stepname=Green] at ([yshift=-4mm]mid.north) {Intermediate state};
\node[dtl] at ([yshift=-8mm]mid.north) {layer outputs\\(reflex, tactical)\\one record per module\\with a vector hash\\feeder provenance};
\node[stepname=Orange] at ([yshift=-4mm]out.north) {Output state};
\node[dtl] at ([yshift=-8mm]out.north) {candidates\\selected option\\ranked, forbidden\\rationale\\vector summary};
\node[panel=Teal, minimum width=56mm, minimum height=15mm] (h) at (44mm,-34mm) {};
\node[stepname=Teal] at ([yshift=2mm]h.center) {proof\_hash};
\node[desc, text width=52mm] at ([yshift=-3mm]h.center) {SHA-256 of the sorted-key JSON\\of every field but itself};
\node[panel=Grey, minimum width=30mm, minimum height=13mm, align=center] (prev) at (-9mm,-34mm) {previous\\proof\_hash};
\draw[flow] (in.south) -- (h.north west);
\draw[flow] (mid.south) -- (h.north);
\draw[flow] (out.south) -- (h.north east);
\draw[flow] (prev.east) -- node[flowlabel, below]{chain} (h.west);
\end{tikzpicture}

\caption{The field groups of a \code{DecisionProof} and how its hash is formed. The previous
proof's hash is one of the hashed fields, so a chain of proofs breaks at the first altered
record.}
\label{fig:deme-proof}
\end{figure}

Let $P$ be the proof's canonical dictionary with the field \code{proof\_hash} removed, and let
$\mathrm{JSON}_{\mathrm{c}}$ be serialization with sorted keys and no whitespace. Then
\begin{equation}
\code{proof\_hash} = \mathrm{SHA256}\bigl(\mathrm{JSON}_{\mathrm{c}}(P)\bigr).
\label{eq:deme-proofhash}
\end{equation}
The canonical dictionary sorts the active module names, the candidates and the forbidden list,
and keeps the ranking in order. \code{verify\_hash} recomputes the hash. \code{DecisionProofChain}
sets each new proof's \code{previous\_proof\_hash} to the last proof's hash before finalizing,
and \code{verify\_chain} checks every hash and every link.
% src: lib:src/erisml/ethics/decision_proof.py:203-271, 386-453

Two properties of the hash matter to an auditor. The pipeline finalizes each proof by itself,
not through a chain, so a pipeline proof has no previous hash unless the caller chains it.
Each proof also gets a random UUID and a wall-clock timestamp, both inside the hashed fields,
and the layer outputs carry measured durations. The same facts therefore give a different
\code{proof\_hash} on every run. The hash shows that a stored record is intact. It does not let
anyone recompute a decision's hash from its inputs.
% src: lib:src/erisml/ethics/decision_proof.py:138-144 (uuid4, now), 249-271
% src: lib:src/erisml/ethics/layers/pipeline.py:307 (proof.finalize, no chain)

\section{Governance, tensor, coalition and uncertainty code}\label{sec:deme-v3}

\paragraph{Governance.} The file \code{governance.py} in the \code{ethics} package is a stub. It
defines a configuration class with a mode, a threshold and weights. Its aggregation function
returns the first judgement, and its selection function returns the first option. Next to the
file is a package directory with the same name. Python imports the package, not the file, and we
confirmed that the module path \code{erisml.ethics.governance} resolves to the package's
\code{\_\_init\_\_.py}. The stub file is unreachable by that name. The package holds the working
versions. Its function \code{aggregate\_moral\_vectors} computes the same weighted mean as
Equation~\eqref{eq:deme-aggregate}, but with a weight looked up per module name, tier and
stakeholder. Its configuration class \code{GovernanceConfigV2} adds dimension weights and lexical
priority layers.
% src: lib:src/erisml/ethics/governance.py:1-38
% src: lib:src/erisml/ethics/governance/__init__.py:1-25
% src: lib:src/erisml/ethics/governance/aggregation_v2.py:57-113
% src: lib:src/erisml/ethics/governance/config_v2.py:8-9, 96-99, 129-158
% src: checked 2026-10-03 with python, erisml.ethics.governance.__file__ ends in governance/__init__.py

\paragraph{Moral tensor.} \code{MoralTensor} generalizes the vector to ranks one to six. The first
axis always has length nine and holds the dimensions of Figure~\ref{fig:deme-vector}. The
default axes are $(k)$, $(k,n)$, $(k,n,\tau)$, $(k,n,a,c)$, $(k,n,\tau,a,c)$ and
$(k,n,\tau,a,c,s)$, where $n$ indexes parties, $\tau$ time, $a$ actions, $c$ coalition
configurations and $s$ Monte Carlo samples. The rank-2 tensor is the per-stakeholder form.
\code{from\_moral\_vectors} stacks one vector per party into a $9 \times n$ array and records
each party's veto flags with that party's column as the veto location. Contraction along an
axis with weights $w$ is a weighted mean,
\begin{equation}
T'_{\dots} = \operatorname{clip}_{[0,1]}\Bigl(\sum_{i} \hat w_i\, T_{\dots i \dots}\Bigr),
\qquad \hat w_i = \frac{w_i}{\sum_j w_j},
\label{eq:deme-contract}
\end{equation}
with uniform weights by default. \code{to\_vector} collapses a tensor to a vector by mean, maximum,
minimum, weighted contraction or a single party. Its docstring calls the minimum the worst case
per dimension. For \code{physical\_harm} the minimum is the least harm, so the minimum strategy
is the most optimistic choice on that dimension and the most pessimistic on the other eight.
% src: lib:src/erisml/ethics/moral_tensor.py:1-25, 79-93 (DEFAULT_AXIS_NAMES)
% src: lib:src/erisml/ethics/moral_tensor.py:183-236 (fields, veto_locations), 477-539 (from_moral_vectors)
% src: lib:src/erisml/ethics/moral_tensor.py:865-938 (contract), 940-1034 (to_vector, "min" docstring at 953)

\paragraph{Coalitions.} \code{coalition.py} works on rank-4 tensors $(k,n,a,c)$ and treats
coalition configuration $0$ as the grand coalition. \code{check\_coalition\_stability} picks one
welfare dimension $d$ and averages it over parties and actions for each configuration,
$W_c = \operatorname{mean}_{n,a} T_{d,n,a,c}$. Configuration $c$ blocks the grand coalition when
its improvement $\Delta_c = W_c - W_0$ (or $W_0 - W_c$ when $d$ is harm) exceeds a threshold,
$0.05$ by default. The stability score is
\begin{equation}
1 - \frac{|B|}{\max(1,\ C-1)},
\label{eq:deme-stability}
\end{equation}
where $B$ is the set of blocking configurations and $C$ the number of configurations. The default
welfare dimension is \code{physical\_harm}. This is a comparison of averages per configuration.
It is not a computation of the core of a cooperative game, although the result sets its
\code{core\_non\_empty} field equal to its stability flag.
% src: lib:src/erisml/ethics/coalition.py:5-23, 453-550

\paragraph{Uncertainty.} \code{uncertainty.py} holds Monte Carlo samples of moral values and
summarizes them by expected value, variance, percentiles, confidence intervals, worst and best
case, and conditional value at risk~\cite{rockafellar2000cvar}. For $S$ samples and tail
probability $\alpha$ the array version computes
\begin{equation}
\mathrm{CVaR}_\alpha = \frac{1}{m}\sum_{i=1}^{m} x_{(i)},\qquad m = \max\bigl(1, \lceil \alpha S \rceil\bigr),
\label{eq:deme-cvar}
\end{equation}
the mean of the $m$ lowest samples, with $\alpha = 0.05$ by default. The module docstring gives
the rank-5 shape as $(d,n,t,c,s)$. That differs from the tensor's rank-5 default
$(k,n,\tau,a,c)$. The architecture document records that the rank-5 conflict was resolved in
favour of the tensor's order, with samples at rank 6. The docstring in \code{uncertainty.py}
still shows the older layout.
% src: lib:src/erisml/ethics/uncertainty.py:5-26, 55-63, 162-179, 477-508
% src: lib:docs/moralvector_v2_architecture.md:85-89 (rank-5 reconciled)

None of the tensor, coalition or uncertainty code runs in the twin. The twin's gate works on
single vectors, one per option.

\section{The twin uses DEME as its output gate}\label{sec:deme-twin}

The twin's robot proposes one action from the scene's allowed set. Before the body acts, the
class \code{OutputGate} in \code{twin/output\_gate.py} puts that proposal and every other allowed
action through a \code{DEMEPipeline}. The pipeline holds three modules,
\code{GenevaEMV2}, \code{AutonomyConsentEMV2}, and \code{TragicConflictEM} wrapped by
\code{V1ToV2Adapter} at tier 3, and its tactical layer is built with \code{fail\_closed} set.
Figure~\ref{fig:deme-twin-gate} shows where the gate's facts come from and the four outcomes it
can produce.
% src: gtc:twin/output_gate.py:1-23 (docstring), 49-63 (_deme)
% src: gtc:twin/brain.py:319-321 (the brain builds the gate), 623 (every decision passes it)

\begin{figure}[t]
\centering
% Chapter 6. The twin's output gate: where the facts come from and the four outcomes, in the order
% OutputGate.check tests them.
% src: gtc-prototype twin/output_gate.py:98-168 (facts), 179-250 (check)
% src: gtc-prototype twin/brain.py:319-321, 623 (the chooser's proposal reaches the gate)
\begin{tikzpicture}[house,
  gbox/.style={minimum width=31mm, minimum height=14mm},
  res/.style={minimum width=33mm, minimum height=32mm},
  dtl/.style={detail, text width=30mm, anchor=south}]
\node[panel=Grey, gbox] (cls) at (0,0) {};
\node[stepname=Grey] at ([yshift=2.2mm]cls.center) {Classifier (model)};
\node[desc] at ([yshift=-2.5mm]cls.center) {proposes events};
\node[panel=Blue, gbox] (st) at (0,-20mm) {};
\node[stepname=Blue] at ([yshift=2.2mm]st.center) {Canonical state};
\node[desc] at ([yshift=-2.5mm]st.center) {norms, machines, events};
\node[panel=Purple, gbox] (sc) at (0,-40mm) {};
\node[stepname=Purple] at ([yshift=2.2mm]sc.center) {Scene profiles};
\node[desc] at ([yshift=-2.5mm]sc.center) {capability\_ethics, refusals};
\node[panel=Grey, gbox] (ch) at (42mm,0) {};
\node[stepname=Grey] at ([yshift=2.2mm]ch.center) {Chooser (model)};
\node[desc] at ([yshift=-2.5mm]ch.center) {one allowed action};
\node[panel=Green, gbox] (f) at (42mm,-30mm) {};
\node[stepname=Green] at ([yshift=2.2mm]f.center) {EthicalFacts};
\node[desc] at ([yshift=-2.5mm]f.center) {one per allowed action};
\node[panel=Teal, gbox, minimum width=38mm, minimum height=30mm] (d) at (90mm,-15mm) {};
\node[stepname=Teal] at ([yshift=10mm]d.center) {DEMEPipeline};
\node[desc, text width=36mm] at ([yshift=-3mm]d.center) {GenevaEMV2 tier 0\\AutonomyConsentEMV2 tier 2\\TragicConflictEM tier 3\\fail\_closed};
\draw[untrusted] (cls) -- node[flowlabel, right]{accepted events} (st);
\draw[flow] (st.east) -- (f.west);
\draw[flow] (sc.east) -- (f.west);
\draw[untrusted] (ch.east) -- node[flowlabel, above]{proposal} (ch.east -| d.west);
\draw[flow] (f.east) -- (f.east -| d.west);
% outcomes
\foreach \c/\x/\n/\t/\b in {
  Red/0/1/DEME raised/{every option vetoed\\wait\_and\_observe\\gate\_error kept},
  Amber/36/2/Tragic and deliberate/{contact the\\monitoring centre\\with the triggers},
  Orange/72/3/Forbidden or not allowed/{best ranked allowed\\action, else\\wait\_and\_observe},
  Green/108/4/Otherwise/{the proposal\\passes unchanged}}{
  \node[panel=\c, res] (o\n) at (\x mm,-74mm) {};
  \node[step=\c, minimum size=5.5mm] at ([yshift=-4.5mm]o\n.north) {\n};
  \node[stepname=\c, text width=31mm, align=center] at ([yshift=-12mm]o\n.north) {\t};
  \node[dtl] at ([yshift=1.5mm]o\n.south) {\b};
}
\draw[draw=hsFlow, line width=1.6pt] (d.south) -- (d.south |- 0,-51mm);
\draw[draw=hsFlow, line width=1.6pt] (o1.north |- 0,-51mm) -- (o4.north |- 0,-51mm);
\foreach \n in {1,...,4}{\draw[flow] (o\n.north |- 0,-51mm) -- (o\n.north);}
\node[flowlabel, anchor=south west] at (d.south |- 0,-51mm) {tested in this order};
\end{tikzpicture}

\caption{The twin's output gate. Code derives the facts from the canonical state and the scene's
declared profiles, and the gate tests four outcomes in order. The two model-written inputs, shown
dashed, reach DEME only as accepted events and as a proposal, and no outcome leaves the allowed
set.}
\label{fig:deme-twin-gate}
\end{figure}

\subsection{The facts are derived, not guessed}

For each allowed action the gate builds one \code{EthicalFacts} record in code, from the
canonical moral state and from the scene's \code{capability\_ethics} block. No model writes
them. Table~\ref{tab:deme-twin-facts} gives every field the gate sets and its source.
% src: gtc:twin/output_gate.py:98-168

\begin{table}[t]
\centering
\footnotesize
\begin{tabularx}{\linewidth}{l X}
\toprule
Field & Value set by \code{OutputGate.facts} \\
\midrule
\code{expected\_benefit} & by the lowest tier of an obligation in force for the action, $1.0, 0.75, 0.5, 0.3$ for tiers $0$ to $3$, else $0.2$ for the scene's default action and $0.1$ otherwise \\
\code{urgency} & equal to the benefit when an obligation is in force, else $0$ \\
\code{expected\_harm} & the \code{harm} declared for the capability, else $0$ \\
\code{violates\_rights} & true when one of Margaret's standing refusals covers the action \\
\code{has\_valid\_consent} & false when the action is refused, or when it touches privacy or her body and neither her consent machine reads \code{obtained} nor the privacy promise has been lifted \\
\code{violates\_explicit\_rule} & true when the action is prohibited in the current state \\
\code{role\_duty\_conflict} & true when an obligation of tier $0$ or $1$ is in force, the action is not obliged and it is not \code{request\_authority} \\
\code{coercion\_or\_undue\_influence} & true when the capability is declared coercive and its \code{justified\_when} conditions do not all hold \\
\code{privacy\_invasion\_level} & $0.8$ for a capability declared private, else $0$ \\
\code{secondary\_use\_without\_consent} & true for a private capability while the privacy promise is active \\
\code{epistemic\_status} & uncertainty $0.2$, evidence quality high, for every action \\
\bottomrule
\end{tabularx}
\caption{The facts the twin's gate builds for each allowed action, and where each value comes
from. Every value is a function of the runtime state and the scene file.}
\label{tab:deme-twin-facts}
\end{table}
% src: gtc:twin/output_gate.py:29 (TIER_BENEFIT), 89-96 (_tier_of), 109-168

The scene declares the profiles. Physical assistance has harm $0.1$ and touches Margaret.
Separating the dog has harm $0.2$ and driving off an animal $0.4$. Restraining a person has harm
$0.4$ and is coercive, justified when restraint is authorized. The stun device and the spray have
harm $0.6$ and are coercive, justified only when restraint is authorized and a severe attack has
been measured. Recording, sharing data and entering the bedroom touch privacy. Unlocking the
medication box has harm $0.2$. Capabilities not listed are treated as benign.
% src: gtc:twin/scene/margaret_home.erisml:1021-1034

The derivation is fixed code, but its inputs are not all fixed. The canonical state includes
events that a language-model classifier proposed and the runtime accepted. A refusal is such an
event. The gate's facts are only as true as the events behind them, which is why the refusal rule
below reads the event history with care.

\subsection{Margaret's refusals bind while she can voice them}

The scene lists seven kinds of refusal and the capabilities each covers. A refusal of emergency
services covers the call to emergency services. Refusals of physical help, of entry to her room,
of recording, of data sharing, of staying inside and of medication cover one capability each.
While a refusal stands, the gate's facts for a covered action record a rights violation without
consent. \code{GenevaEMV2} then vetoes the action at tier 0.
% src: gtc:twin/scene/margaret_home.erisml:1003-1020
% src: gtc:twin/output_gate.py:17-19, 145-147

\begin{figure}[t]
\centering
% Chapter 6. How OutputGate.refused() reads Margaret's refusals from the event history.
% src: gtc-prototype twin/output_gate.py:65-87
% src: gtc-prototype twin/scene/margaret_home.erisml:1003-1020 (refusals, refusal_lapses_on)
\begin{tikzpicture}[house,
  st/.style={minimum width=36mm, minimum height=17mm},
  lbl/.style={flowlabel, align=center}]
\node[panel=Green, st] (a) at (0,0) {};
\node[stepname=Green] at ([yshift=3mm]a.center) {She can voice};
\node[desc] at ([yshift=-2.5mm]a.center) {no refusal in force};
\node[panel=Orange, st] (b) at (62mm,0) {};
\node[stepname=Orange] at ([yshift=3mm]b.center) {Refusal stands};
\node[desc, text width=33mm] at ([yshift=-3mm]b.center) {its capabilities are\\facts of a rights violation};
\node[panel=Grey, st] (c) at (31mm,-36mm) {};
\node[stepname=Grey] at ([yshift=3mm]c.center) {Lapsed};
\node[desc, text width=33mm] at ([yshift=-3mm]c.center) {every refusal cleared\\new refusals ignored};
\draw[untrusted] (a) -- node[lbl, above]{refusal\_made} (b);
\draw[flow] (b.south) -- node[lbl, right]{unresponsive or\\check\_in\_unanswered} (c.east);
\draw[flow] ([xshift=8mm]a.south) -- node[lbl, right, pos=0.35]{lapse} (c.north west);
\draw[flow] (c.west) -| node[lbl, left, pos=0.75]{check\_in\_answered} ([xshift=-10mm]a.south);
\draw[untrusted] (c.south east) .. controls +(14mm,-12mm) and +(-14mm,-12mm) .. node[lbl, below]{refusal\_made\\(no effect)} (c.south west);
\node[detail, text width=40mm, anchor=south] at ([yshift=2mm]b.north) {covers nothing while a hazard\\in its not\_during list holds};
\end{tikzpicture}

\caption{How the gate decides which refusals stand. A lapse event clears every refusal, and only
an answered check-in lets a new refusal count again. The refusal events are model-classified,
shown dashed, which is why a lapse must be followed by evidence that she can speak.}
\label{fig:deme-refusal}
\end{figure}

The method \code{refused} replays the event history and keeps a flag for whether Margaret can
voice a refusal (Figure~\ref{fig:deme-refusal}). A lapse event, \code{unresponsive} or
\code{check\_in\_unanswered} in the scene, clears every refusal and sets the flag false. The
event \code{check\_in\_answered} sets it true. A \code{refusal\_made} event by Margaret counts
only while the flag is true. A refusal also covers nothing while a hazard named in its
\code{not\_during} list holds. The scene uses this for emergency services, whose refusal does not
cover a call during danger in the home or an attack. The scene's comment records that refusing an
ambulance does not refuse the fire service in a fire, and that lifting a refusal grants nothing,
since emergency services still need the governor's ruling.
% src: gtc:twin/output_gate.py:65-87
% src: gtc:twin/scene/margaret_home.erisml:1003-1013, 1020

The re-arming rule came from a failure. In development run dev12b, case d27, the classifier
re-read Margaret's words ``you don't need to call anyone'', said before she collapsed, as a fresh
refusal while she lay unresponsive, and DEME vetoed the authorized emergency call. Under the
current rule a refusal recorded after a lapse has no effect until she answers a check-in. Three
tests fix the behaviour. A refusal she can no longer voice does not bind. Her old words re-read
after a lapse do not refuse again. A refusal made after an answered check-in binds.
% src: gtc:twin/output_gate.py:66-71 (docstring citing dev12b d27)
% src: gtc:tests/test_twin_ladder.py:1135-1158

\subsection{Tragic choices go to the monitoring centre}

For every check the gate also computes the tragic-conflict index $\tau$ of the proposal. When the
decision is deliberate, $\tau \ge 0.55$, the proposal is not itself the call to the monitoring
centre, and that call is allowed and not forbidden, the gate replaces the proposal with
\code{contact\_monitoring\_center}. The message names the proposal and its triggers. This
happens whether or not DEME vetoed the proposal.
% src: gtc:twin/output_gate.py:170-177 (tragic), 199-225 (routing)

The emergency call against a standing refusal is the case the tests use. The scene's obligations
to call emergency services carry priority tier 0, so the gate sets benefit and urgency to $1.0$.
The refusal sets the rights violation and removes consent. No harm is declared for the call.
By Equation~\eqref{eq:deme-tragic},
\begin{equation}
\tau = 0.20 + 0.25 + 0.10 = 0.55,
\end{equation}
which meets the threshold. The test asserts the high flag, the routing and that the message
mentions a tragic conflict. The arithmetic here is ours, from the code and the scene. The test
records only the outcome.
% src: gtc:twin/scene/margaret_home.erisml:190-206 (EMS obligations at priority_tier 0)
% src: gtc:tests/test_twin_ladder.py:1118-1132

\subsection{Reflexes are gated but not redirected}

The twin's reflexes respond to a measured attack on Margaret without a model. The brain passes
each reflex action through the same gate with \code{deliberate} false. DEME may still veto a
reflex action and replace it, but the gate never routes a reflex to the monitoring centre,
because a reflex must stay fast. A test checks that a reflex emergency call against a standing
refusal is vetoed, replaced by an allowed action and not routed.
% src: gtc:twin/brain.py:472-473
% src: gtc:twin/output_gate.py:21, 179-206
% src: gtc:tests/test_twin_ladder.py:1161-1172

\subsection{Replacement stays inside the allowed set}

If DEME forbids the proposal, or the proposal is not in the allowed set, the gate takes the first
action in DEME's ranking that is neither forbidden nor outside the allowed set. If there is none
it takes \code{wait\_and\_observe} when that is allowed, else the scene's default action. The gate
does not use the pipeline's \code{selected\_option\_id}, so the mismatch noted in
Section~\ref{sec:deme-pipeline} cannot reach the body. The gate record keeps the proposal, the
veto, the forbidden list, the first six ranked options, the tragic index, any module failures
and the proof hash. The Lean model of the authority path proves that the gate's output is always
permitted by the scene, whatever DEME vetoes and however it ranks.
% src: gtc:twin/output_gate.py:226-250
% src: gtc:formal/twin-containment/README.md:16 (gate_permitted and companions)
% src: gtc:formal/twin-containment/TwinContainment.lean:379

\subsection{The gate fails closed at three levels}

A gate that cannot judge must not look like a gate that approves. The twin closes three ways.
Inside the tactical layer, \code{fail\_closed} turns a failing module into a veto of that option.
Around the pipeline, any exception from DEME makes the gate return the idle action with every
option recorded as vetoed and the error text kept in \code{gate\_error}. Before serving, the
brain runs a self-test that puts the default action through the gate on the scene's reset state
and refuses to start if the record shows a gate error or any module failure. Tests cover each
level.
% src: gtc:twin/output_gate.py:58-62 (fail_closed=True), 189-198 (exception path)
% src: gtc:twin/brain.py:579-588 (self_test); gtc:twin/service.py:385
% src: gtc:tests/test_twin_ladder.py:1181-1185, 1205-1230

The self-test exists because of an incident. Its docstring records that development run dev8 on
2026-10-02 ran a whole development set on a gate that raised on every call, against a stale
copy of the library, and nothing reported it.
% src: gtc:twin/brain.py:579-582

\section{Which parts are deterministic}\label{sec:deme-deterministic}

The prior-art review for this work asks that the report state which DEME modules are
deterministic, because a veto gate whose modules are models falls in the same class as model
guardrails~\cite{rebedea2023nemo,inan2023llamaguard}. In the twin the answer is all of them.
% src: gtc:paper/aies2027/prior_art.md:103

\begin{itemize}
\item \code{GenevaEMV2}, \code{AutonomyConsentEMV2} and \code{TragicConflictEM} are fixed
arithmetic on the facts, as the equations of Section~\ref{sec:deme-modules} show. None calls a
model.
\item The reflex layer, the deontic gate, the tactical aggregation, the ranking and the
baseline rule are fixed code over the facts and the module outputs.
\item The twin's facts are computed by \code{OutputGate.facts} from the runtime state and the
scene file. No model writes them.
\item The refusal replay, the tragic routing and the replacement rule are fixed code.
\end{itemize}
% src: lib:src/erisml/ethics/modules/tier0/geneva_em.py:73-186
% src: lib:src/erisml/ethics/modules/tier2/autonomy_consent_em.py:56-133
% src: lib:src/erisml/ethics/modules/greek_tragedy_tragic_conflict_em.py:63-142
% src: gtc:twin/output_gate.py:65-250

Three things are not deterministic in the same sense. The decision proof's hash changes on every
run, because of its random identifier, its timestamp and its measured durations, while the
forbidden set and the ranking do not. The canonical state that feeds the facts includes events
proposed by a language-model classifier, so a misclassified event becomes a wrong fact, as case
d27 showed. And the library supports dimensions scored by learned feeders, which the twin does
not use. A deployment that enabled them would put a model inside the vector, and the proof's
feeder record exists to say which encoder, at which checkpoint, passed which bar.
% src: lib:src/erisml/ethics/decision_proof.py:80-122, 138-144
% src: gtc:twin/output_gate.py:66-71

The library's own default leaves the tactical layer fail-open. The twin overrides it. Another
caller of \code{DEMEPipeline} gets the earlier behaviour, in which a module that raises is
recorded and logged but the option is judged by the modules that answered, unless that caller
sets \code{fail\_closed}.
% src: lib:src/erisml/ethics/layers/tactical.py:68-76
